\documentclass[10pt,a4paper]{amsart}
\usepackage[margin=19mm]{geometry}
\usepackage{amsmath,amssymb,mathtools}
\usepackage[scr=rsfs,cal=euler]{mathalfa}
\usepackage{xfrac}
\usepackage[unicode,hidelinks,bookmarks=false]{hyperref}
\newcommand{\ie}{i.e.\ }
\newcommand{\cf}{cf.\ }
\newcommand{\resp}{respectively\ }
\newcommand{\Subsection}{\subsection}
\providecommand{\nobreakdash}{\nobreak\hbox{-}}
\setlength{\emergencystretch}{2em}
\multlinegap0pt
\usepackage{dsfont}
\usepackage{xspace}
\usepackage{enumitem}
\usepackage{thmtools}
\usepackage{thm-restate}
\theoremstyle{plain}
\newtheorem{theorem}{Theorem}
\newtheorem{lemma}{Lemma}
\newtheorem{proposition}{Proposition}
\theoremstyle{definition}
\newtheorem{definition}{Definition}
\theoremstyle{remark}
\newtheorem{remark}{Remark}
\usepackage[capitalize]{cleveref}
\newcommand{\N}{\mathbb{N}}
\newcommand{\R}{\mathbb{R}}
\newcommand{\E}{\mathbb{E}}
\newcommand{\pr}{\mathbb{P}}
\renewcommand{\Pr}{\pr}
\newcommand{\ind}[1]{\mathds{1}\left\{#1\right\}}
\newcommand{\Exp}[1]{\mathbb{E} \left[ #1 \right]}
\newcommand{\BDist}{\mathcal{B}}
\newcommand{\NDist}{\mathcal{N}}
\newcommand{\Bmean}{\mu_{\BDist}}
\newcommand{\Nmean}{\mu_{\NDist}}
\newcommand{\Bstd}{\sigma_{\BDist}}
\newcommand{\Nmax}{N_{\max}}
\newcommand{\Bmax}{B_{\max}}
\newcommand{\maxAlloc}{A_{\max}}
\newcommand{\Deff}{\Delta_{\text{Efficiency}}\xspace}
\let\oldparagraph\paragraph
\renewcommand{\paragraph}[1]{\subsubsection*{#1}}
\title[Policy bounds for finite expected hitting times and boundary hit probabilities]{Sequential Fair Allocation With Replenishments: A Little Envy Goes An Exponentially Long Way: the policy-constants lemma (finite expected hitting times and boundary hit probabilities)}
\author{Chido Onyeze, Sean R. Sinclair, Chamsi Hssaine, Siddhartha Banerjee}
\date{Source: arXiv:2508.21753v1 (CC BY 4.0)}
\begin{document}
\begin{abstract}
We study the trade-off between envy and inefficiency in repeated resource allocation settings with stochastic replenishments, motivated by real-world systems such as food banks and medical supply chains. Specifically, we consider a model in which a decision-maker faced with stochastic demand and resource donations must trade off between an {\it equitable} and {\it efficient} allocation of resources over an infinite horizon. The decision-maker has access to storage with fixed capacity $M$, and incurs efficiency losses when storage is empty (stockouts) or full (overflows). We provide a nearly tight (up to constant factors) characterization of achievable envy-inefficiency pairs. Namely, we introduce a class of Bang-Bang control policies whose inefficiency exhibits a sharp phase transition, dropping from $\Theta(1/M)$ when $\Delta = 0$ to $e^{-\Omega(\Delta M)}$ when $\Delta > 0$, where $\Delta$ is used to denote the target envy of the policy. We complement this with matching lower bounds, demonstrating that the trade-off is driven by {\it supply}, as opposed to {\it demand} uncertainty. Our results demonstrate that envy-inefficiency trade-offs not only persist in settings with dynamic replenishment, but are shaped by the decision-maker's available capacity, and are therefore qualitatively different compared to previously studied settings with fixed supply.


\end{abstract}
\maketitle

\noindent\emph{Source and licence.} Sequential Fair Allocation With Replenishments: A Little Envy Goes An Exponentially Long Way. Version arXiv:2508.21753v1 (CC BY 4.0), distributed under the Creative Commons Attribution 4.0 International licence, http://creativecommons.org/licenses/by/4.0/, as recorded in the arXiv licence metadata for this exact version and shown on the abstract page https://arxiv.org/abs/2508.21753v1. Full rights notice: licenses/arxiv-2508.21753-CC-BY-4.0.txt.\par\smallskip

\noindent\textit{Editorial scope.} Counted unit: the policy-constants lemma of the paper, printed in the arXiv version as Lemma 3 (source0.tex lines 1413--1426, source label \texttt{lem:policy\_constants}, restatable name PolicyBounds, in Appendix B.1), delivered together with the paper's complete original proof of it (source0.tex lines 1427--1491, one \texttt{proof} environment of 65 lines with no gap and no deferral). No mathematics is written, repaired or re-derived: statement and proof are the paper's own text line for line, in the paper's own notation ($M$, $S_t$, $B_t$, $N_t$, $A_t$, $Z_t(\alpha)$, $\tau$, $\tau_M$, $\tau_0$, $E(S)$, $p_M$, $p_0$, $Q_t$, $\epsilon$, $\delta$). The paper states the lemma inside a \texttt{restatable} environment of the \texttt{thm-restate} package and never calls \texttt{\textbackslash PolicyBounds*} anywhere in the source, so the retained environment is the lemma's only printed instance and the wrapper reproduces that environment verbatim. Required context retained uncounted: the whole of the paper's Preliminary section heading and its model primitives (source0.tex lines 407--408 and 412: the storage capacity $M$, the initial level $S_0$, the donated budget $B_t$ and the arriving agents $N_t$), the inventory-update relation with the drift $Z_t$ and the bounds $Z_{\min}=\maxAlloc\Nmax$, $Z_{\max}=\Bmax$ (lines 414--418), the i.i.d.\ supply and demand distributions and their means (lines 420--421), the definition of an allocation policy and of a time-homogeneous policy (lines 427--430), the overflow and stockout quantities $W_t$ and $V_t$ (lines 446--451) and the paper's definition of inefficiency with the long-run averages $\overline{W}$ and $\overline{V}$ (lines 452--464), the subsection heading and label under which the paper introduces the hitting times together with the definitions of $H_M$, $H_0$ and of the per-round drift $Z_t(\alpha)$ (lines 615--619), the paper's own itemised definitions of $\tau(S)$, $E(S)$, $\tau_M(S)$, $\tau_0(S)$, $p_M$ and $p_0$ (lines 631--635), the paper's sentence introducing the counted lemma (line 1411), and the statement of the paper's random-walk-equivalence lemma (lines 726--729) that fixes the unreflected process $Q_t=\sum_{s=1}^t Z_s+S_0$ and the coupling $E(S)=\E[\tau\mid S_0=S]$ used in part (i) of the counted proof; the paper's abstract is reproduced verbatim as front matter (line 321). Every definition, numbered display, label and structural statement of those lines is retained, unchanged. Omitted from this package and not redistributed: the paper's title-page block and author affiliations, its keywords, its bibliography and the external \texttt{main.bbl} that ships with the e-print, its introduction and related-work sections (source0.tex lines 341--406), the warm-up section (lines 480--606), the remainder of the achievability section (lines 609--614, 620--630, 636--725, 730--980), the lower-bound, extension, simulation and conclusion sections (lines 981--1263), the appendix material other than the counted lemma, its proof and the paper's sentence introducing it (lines 1305--1410, 1412 and 1492--2077), and the paper's closing rules. Within the retained region the following passages are also omitted: the preliminary section's opening sentence, which points forward to the paper's extension section (lines 409--411), the paragraph fixing the per-agent allocation choice together with its footnote citation (line 413), the remark on the i.i.d.\ assumption (lines 423--425), the ex-post envy definition and its motivation, which the counted lemma does not use (lines 431--445), and the achievability section's opening paragraph (lines 609--614), which points forward to material not retained. Nothing else of the counted statement's or its proof's dependencies is omitted; the counted lemma is a statement about the inventory process alone and uses no envy quantity. Reference adaptation: none was needed and none was made. Every reference inside the delivered unit resolves inside it: the label \texttt{eq:state\_dynamics} of the retained inventory-update display, which the same retained paragraph cites; the label \texttt{ssec:timehomogen} of the retained subsection heading, which the paper's retained sentence at line 1411 cites; and the label \texttt{eq:expect-tau0} of the numbered display inside the counted proof, which the proof's own later sentence cites. The paper's own \texttt{cleveref} package and its \texttt{\textbackslash Cref} and \texttt{\textbackslash cref} commands are reproduced as the paper uses them, so those cross-references print with the paper's own capitalisation convention. Labels that are retained but not cited anywhere in the unit (\texttt{sec:model}, \texttt{eq:waste}, \texttt{eq:subsidy}, \texttt{def:efficiency}, \texttt{eq:deff}, \texttt{lem:policy\_constants}, \texttt{lem:Random\_walk\_equivalence\_1}) are the paper's own anchors and produce no unresolved reference. No unresolved reference or citation is produced. Citations: the retained spans contain no citation command, so the delivered unit prints no bibliography and no bibliography entry is transcribed or redistributed; the paper's own 44-entry \texttt{main.bbl} and its \texttt{references.bib} are not copied into the package. Printed slips kept literal: no printed word, symbol, hypothesis or number of the counted statement or of its proof was altered. In particular item (iii) of the statement is delivered exactly as printed, ``$p_0 \in (0,1)$ and $p_M (0,1)$'', with the missing membership symbol that the source itself prints there; the proof's own mixture of notations is delivered as printed ($\Exp{\cdot}$ in the statement and in the first display, $\mathbb{E}[\cdot]$ and $\mathbb{E}\big[\cdot\big]$ later in the proof); the paper's own line break inside the definition of the event $E^-_L$ is kept; and the paper's own renaming of the expected hitting time ($E(S)$ for $\Exp{\tau(S)}$, and $E(M)$, $E(0)$ for the boundary values) is kept. Layout and class: the paper's own class is \texttt{informs3} with the options \texttt{mnsc} and \texttt{nonblindrev}, and the e-print ships \texttt{informs3.cls} together with \texttt{informs2014.bst}; neither the class file nor the style file is redistributed with this package and neither is used to build it. The wrapper is the lane's pinned \texttt{amsart} editorial wrapper at 10pt on A4, to which this package adds the paper's own macro block (its definitions of $\N$, $\R$, $\E$, $\pr$ with the paper's own renewal of \texttt{\textbackslash Pr}, the indicator \texttt{\textbackslash ind} through \texttt{dsfont}, the expectation \texttt{\textbackslash Exp}, the supply and demand distributions \texttt{\textbackslash BDist} and \texttt{\textbackslash NDist} with their means, standard deviations and upper bounds, the maximum allocation, \texttt{\textbackslash Deff}, and the paper's own renewal of \texttt{\textbackslash paragraph} as an unnumbered run-in heading), the paper's own \texttt{thmtools} and \texttt{thm-restate} loading, its own \texttt{cleveref} loading and its \texttt{enumitem} lists, and declares the theorem environments the retained statements need. Because the paper's class supplies those environments, the wrapper declares them itself; the delivered numbering is therefore the unit's own (the two retained lemmas print as Lemma 1 and Lemma 2 in the order delivered, where the paper prints the random-walk-equivalence lemma as Lemma 1 and the counted lemma as Lemma 3, and the retained inefficiency definition prints as Definition 1 of the unit). The paper's fonts (its \texttt{tgtermes}, \texttt{newtxtext} and \texttt{newtxmath} families) are replaced by the wrapper's Latin Modern, which is a page-appearance difference of the wrapper only; no formula, symbol, hypothesis, constant or argument changes. Assets: the paper has nine graphics-inclusion commands (one in its introduction, six in its simulations section and two in its appendix simulations) and nine figure environments, four of which are in-source TikZ pictures; none of those eighteen items lies in a retained span, and no retained span contains a graphics-inclusion command, a PDF-page inclusion, an \texttt{input} directive or an \texttt{include} directive. The package declares no asset dependency and no image file is copied into or redistributed with it; all mathematics of the unit is native text with no image dependency. Numbering and cross-references: the delivered unit's numbers are its own throughout. The retained inventory-update display prints as equation (1) in the unit exactly as it does in the article, while the retained overflow, stockout and inefficiency displays print as (2), (3) and (4) in the unit where the article prints them as (3), (4) and (5); the retained numbered display inside the counted proof prints as (5) in the unit where the article prints it as (20); and the retained subsection cross-reference prints as Section 1.1 in the unit where the article prints Section 4.1. The mathematics, hypotheses and wording of those displays and of that cross-reference are unchanged.Licence: version arXiv:2508.21753v1 of this article is distributed under the Creative Commons Attribution 4.0 International licence, http://creativecommons.org/licenses/by/4.0/, as recorded in the arXiv licence metadata for this exact version and shown on the abstract page https://arxiv.org/abs/2508.21753v1. A full rights notice is delivered at licenses/arxiv-2508.21753-CC-BY-4.0.txt. Characters: the retained spans contain no character outside ASCII; the five non-ASCII characters of the source file (three en dashes and two right single quotation marks) all lie outside the retained spans, so no encoding substitution, glyph replacement or missing-character risk arises.
\section{Preliminaries}
\label{sec:model}

\paragraph{Model Primitives.} We consider a decision-maker who manages a warehouse (also referred to as a {\em store}) of a single divisible resource over an infinite horizon of discrete time periods.  The store has a capacity $M > 0$ and initially contains $S_0 \in [0,M]$ units of the resource. At the start of period $t \in \N$, the decision-maker receives an additional budget of $B_t \in \R_{\geq 0}$ resources (also referred to as a {\em donation}); moreover, $N_t\in \N$ homogeneous agents arrive, each requesting a share of the resource.

\begin{equation}
\label{eq:state_dynamics}
S_{t} = (S_{t-1} + B_t - N_t A_t)\Big|_0^{M},
\end{equation}
where we use the notation $x|_0^{M} = \max\{0,\min\{x,M\}\}$ to denote the projection of $x$ in $[0,M]$. For ease of notation, we define the {\it inventory drift} in round $t$ to be $Z_t = B_t-N_tA_t$, noting that $S_{t} = \left.(S_{t-1} + Z_t)\right|_0^{M}$, by \Cref{eq:state_dynamics}. For ease of notation, we let $Z_{\min} = \maxAlloc\Nmax$ and $Z_{\max} = \Bmax$, noting that $Z_t \in [-Z_{\min},Z_{\max}]$ for all $t$.


We assume $B_t$ is drawn i.i.d. from a distribution $\BDist$ supported on a subset of $[0, B_{\max}]$, and let $\Bmean {> 0}$ denote the mean of $\BDist$, known to the decision-maker. Similarly, $N_t$ is drawn i.i.d. from a distribution $\NDist$ supported on a subset of $\{0,\ldots,N_{\max}\}$, and let $\Nmean {> 0}$ denote its known mean.\footnote{The assumption that $N_t$ is discrete and both $N_t$ and $B_t$ are bounded is for technical simplicity. Our results extend naturally for any $\NDist$ and $\BDist$ which are sub-Gaussian.} Finally, we assume $\BDist$ and $\NDist$ are independent of $M$, and that $\BDist$ has strictly positive variance $\Bstd^2$.

\paragraph{Policies.} We define an allocation policy $\mathcal{A}$ to be a mapping from the current round $t$, the realized supply and demand $B_t$ and $N_t$, and the current inventory level $S_{t-1}$, to a per-individual allocation quantity $A_t$.
Policy $\mathcal{A}$ is said to be {\it time-homogeneous} if 
it only depends on the current round $t$ through the inventory level $S_{t-1}$ and the realized demand and supply. 
Our achievability results all involve time-homogeneous policies, while our impossibilities hold under arbitrary policies.

In particular, we define the {\it overflow} (or {\it waste}) $W_t$ incurred in round $t$ as the amount of resources lost due to the inability to store more resources than capacity $M$. Formally:
\begin{equation}\label{eq:waste}W_t = (S_{t-1} + B_t - N_t A_t - M)^+,
\end{equation}
where $(\cdot)^+ = \max\{\cdot, 0\}$ is used to define the positive part. We moreover define the {\it stockout quantity} $V_t$ in round $t$ to be the additional resource the decision-maker needs to acquire in order to allocate $A_t$. Formally:
\begin{equation}\label{eq:subsidy}V_t = (S_{t-1}+B_t-N_t A_t)^-,\end{equation}
where $(\cdot)^- = -\min\{\cdot, 0\}$.
With these definitions in hand, we can define our measure of inefficiency.
\begin{definition}[Inefficiency]
\label{def:efficiency}
Given per-unit overflow cost $h \in \mathbb{R}_{> 0}$ and per-unit stockout cost $b \in \mathbb{R}_{> 0}$, the \emph{inefficiency} of a sequence of allocation decisions $(A_t, t\in \N)$ is given by:
\begin{equation}
\label{eq:deff}
\Deff = \limsup_{T \rightarrow \infty} \left( \frac{1}{T} \sum_{t=1}^T hW_t + bV_t \right) =  h\overline{W} + b\overline{V},
\end{equation}
where $\overline{W}$ and $\overline{V}$ respectively denote the long-run average overflow and stockouts under $(A_t, t \in \mathbb{N})$, i.e.:
\begin{equation*}
\overline{W} = \limsup_{T \rightarrow \infty} \frac{1}{T} \sum_{t=1}^T W_t, \quad \overline{V} = \limsup_{T \rightarrow \infty} \frac{1}{T} \sum_{t=1}^T V_t.
\end{equation*}
\end{definition}

\subsection{Bounds on Overflow and Stockout}
\label{ssec:timehomogen}


We first establish the intuitive fact that the long-run average waste and stockout quantity respectively grow with the fraction of time the warehouse is full (i.e., $S_t = M$) and empty (i.e., $S_t = 0$). For ease of notation, we let $H_M = \lim_{T \rightarrow \infty}\frac{1}{T}\sum_{t=1}^T \ind{S_{t} = M}$ and $H_0 = \lim_{T \rightarrow \infty}\frac{1}{T}\sum_{t=1}^T \ind{S_{t} = 0}$ under any policy $\mathcal{A}$. Additionally, for any $\alpha \in [0,A_{\max}]$, we let $Z_t(\alpha) = B_t - \alpha N_t$ be the drift in round $t$ for any policy that allocates $\alpha$ in that round.

Toward this goal, we introduce the following notation:
\begin{itemize}
    \item For any initial state $S \in [0, M]$, we let $\tau(S) = \inf\left\{t > 0 : S_{t} \in \{0, M\} \mid S_0 = S \right\}$ be the first time $S_t$ hits $0$ or $M$, and define $E(S) = \mathbb{E}[\tau(S)]$. Additionally, let {$\tau_M(S) = \inf\{t > 0 : S_t = M \mid S_0 = S\}$ and $\tau_0(S) = \inf\{t > 0 : S_t = 0 \mid S_0 = S\}$ respectively denote the first time the process reaches $M$ or $0$, starting from state $S$.}
    \item We let $p_M = \pr(S_{\tau(M)} = M \mid S_0 = M)$ be the probability that, given initial state $M$, $S_t$ reaches $M$ before 0. Similarly  $p_0 = \pr(S_{\tau(0)} = 0 \mid S_0 = 0)$ is used to denote the probability that $S_t$ reaches 0 before $M$, given that initial state 0.
\end{itemize}

\begin{restatable}{lemma}{RandomWalkEquivalence}
\label{lem:Random_walk_equivalence_1}
Let $Q_t = \sum_{s=1}^t Z_s + S_0$, and define stopping time \mbox{$\tau = \inf\{t > 0: Q_t \not\in (0, M)\}$}. Then, for any $S \in [0, M]$, $E(S) = \E[\tau \mid S_0 = S]$. Furthermore, \mbox{$p_M = \pr(Q_{\tau} \geq M \mid S_0 = M)$} and $p_0 = \pr(Q_{\tau} \leq 0 \mid S_0 = 0)$.
\end{restatable}

The following lemma provides conditions on the budget and demand distributions $B_t$ and $N_t$ to ensure that the stopping times defined in \cref{ssec:timehomogen} are well-defined and finite.

\begin{restatable}{lemma}{PolicyBounds}
    \label{lem:policy_constants}
    Fix an allocation policy $\mathcal{A}$.  Suppose there exist strictly positive constants $\epsilon$ and $\delta$ such that, for any allocation $\alpha$ made by the policy at time $t$:
    \[
    \Pr(Z_t(\alpha) \geq \epsilon) \geq \delta \quad \text{ and } \quad \Pr(Z_t(\alpha) \leq - \epsilon) \geq \delta.
    \]
    Then, the following properties hold:
    \begin{enumerate}[label=(\roman*)]
    \item  $E(S) < \infty$ for any $S \in [0,M]$,
    \item $\tau_M(S) < \infty$ and $\tau_0(S) < \infty$, almost surely, for any $S \in [0,M]$,
    \item $p_0 \in (0,1)$ and $p_M (0,1)$, and
    \item $\Exp{Z_t(\alpha)^+} \geq \epsilon \delta$ and $\Exp{Z_t(\alpha)^-} \geq \epsilon \delta$.
    \end{enumerate}
\end{restatable}

\begin{proof}
\noindent\textbf{(i) Finite expected hitting time.} Recall, $\tau_M(S)$ and $\tau_0(S)$ respectively denote the first time the process reaches $M$ or $0$, starting from state $S$. To show that $E(S) < \infty$ it suffices to show that $E_0(S) \coloneqq \Exp{\tau_0(S)} < \infty$ for any $S \in [0,M]$, since $\tau(S) \leq \min\{\tau_M(S), \tau_0(S)\}$. In the remainder of the proof, we omit dependence of all quantities on the initial state $S$, for notational simplicity.

By assumption, there exist $\epsilon$ and $\delta$ such that $\Pr(Z_t(\alpha) \leq -\epsilon) \geq \delta$.  Let $L=\lceil M/\epsilon\rceil$.
By the tail-sum formula,
\begin{align*}
\E[\tau_0(S)] = \sum_{k = 0}^\infty \sum_{l=0}^{L-1} \Pr\big(\tau_0(S)> kL + l\big) \leq \sum_{k = 0}^\infty \sum_{l=0}^{L-1} \Pr\big(\tau_0(S)> kL\big) &= L\sum_{k=0}^{\infty}\mathbb{P}\left(\tau_0(S) > kL\right).
\end{align*}
By the law of total probability, and the fact that $\tau_0(S) \leq (k-1)L \implies \tau_0(S) \leq kL$, we have:
\begin{align}\label{eq:expect-tau0}
\mathbb{P}\left(\tau_0(S) \leq kL\right) &= \mathbb{P}\left(\tau_0(S) \leq (k-1)L\right) + \mathbb{P}\left(\tau_0(S) \leq kL , \tau_0(S) > (k-1)L\right)
\end{align}

Consider the event $
E^-_L \;:=\; \{Z_1\le -\epsilon,\; Z_2\le -\epsilon,\ldots, Z_L\le -\epsilon\}.$ Conditioning the second summand on $E_L^-$, we have:
\begin{align*}
\mathbb{P}\left(\tau_0(S) \leq kL\right) &\geq \mathbb{P}\left(\tau_0(S) \leq (k-1)L\right) + \mathbb{P}\left(\tau_0(S) \leq kL , \tau_0(S) > (k-1)L \mid E_L^-\right)\mathbb{P}(E_L^-)
\end{align*}

Note that, if $\tau_0(S) > (k-1)L$, by definition $S_t$ has not reached 0 by time $(k-1)L$. However, on the event $E_L^-$, the unreflected process $Q_t$ decreases by at least $L\epsilon \geq M$. Therefore, it must be the case that $\tau_0(S) \leq kL$, independent of the starting state at $(k-1)L$. Moreover, since $(B_t,N_t)$ are independent across time,
\(\Pr(E^-_L)\ge \delta^L\). Putting these two facts together, we obtain:
\begin{align*}
\mathbb{P}\left(\tau_0(S) \leq kL\right) &\geq \mathbb{P}\left(\tau_0(S) \leq (k-1)L\right) + \mathbb{P}\left( \tau_0(S) > (k-1)L \mid E_L^-\right)\delta^L\\
&= 1-(1-\delta^L)\mathbb{P}\left(\tau_0(S)>(k-1)L\right)
\\
\implies \mathbb{P}\left(\tau_0(S) > kL\right) &\leq (1-\delta^L)\mathbb{P}\left(\tau_0(S) > (k-1)L\right) \leq (1-\delta^L)^k,
\end{align*}
where the final inequality follows from the fact that $\mathbb{P}(\tau_0(S) > 0) = 1$.

Plugging this back into \eqref{eq:expect-tau0}, we obtain:
\begin{align*}
\mathbb{E}[\tau_0(S)] \leq L\sum_{k=0}^{\infty}(1-\delta^L)^k = \frac{L}{\delta^L} < \infty.
\end{align*}

An analogous argument with $E^+_L := \{Z_1\ge \epsilon,\ldots,Z_L\ge \epsilon\}$ gives a finite bound for the expected time to hit $M$, i.e. $\Exp{\tau_M(S)} < \infty$.


\medskip
\noindent\textbf{(ii) Finite hitting times almost surely.}
The previous argument established that $\Exp{\tau_M(S)} < \infty$ and $\Exp{\tau_0(S)} < \infty$. This immediately implies that $\tau_M(S)$ and $\tau_0(S)$ are finite almost surely.

\medskip
\noindent\textbf{(iii) Boundary hit probabilities.} We show $p_M\in(0,1)$; the proof for $p_0$ is symmetric.
By definition,
\[
p_M \;=\; \Pr(S_{\tau(M)}=M \,\big|\, S_0 = M).
\]
A one-step return to $M$ occurs whenever $Z_1\ge \epsilon$, since $M+Z_1 > M \implies S_1 := (M+Z_1)\vert_0^M=M$. Hence, $\tau(M) = 1$, with $S_{\tau(M)} = M$.
Therefore, $p_M \ge \Pr(Z_1\ge \epsilon\mid S_0=M)\ge \delta>0$.

To see that $p_M<1$, note that under event $E^-_L$, the process hits 0 in the first $L$ periods (since the cumulative drift is at most $-L\epsilon \leq -M$), before being able to return to $M$. Moreover, this occurs with probability at least
$\Pr(E^-_L\mid S_0=M)\ge \delta^L$.
Therefore $\Pr(S_{\tau(M)}=0\mid S_0=M)\ge \delta^L>0$, which implies $p_M\le 1-\delta^L<1$.
A symmetric argument starting from $S_0=0$ gives $p_0\in(0,1)$.

\medskip
\noindent\textbf{(iv) Positivity of $\Exp{Z_t(\alpha)^+}$ and $\Exp{Z_t(\alpha)^-}$.}
By assumption:
\[
\mathbb{E}\big[ Z(\alpha)^+ \big] 
\ \ge\ \epsilon \cdot \Pr(Z(\alpha) \geq \epsilon) 
\ \ge\ \epsilon \delta.
\]
An analogous argument gives $\mathbb{E}\big[ Z(\alpha)^- \big] \ge \epsilon\delta$.  
\end{proof}

\end{document}
